\section{The space and its compact-open pieces}
\label{sec:geometry}

Fix \(r\geq1\). Choose discrete sets \(A_i\) with
\(|A_i|=\aleph_i\), for \(0\leq i\leq r\), and set

\[
D_i=A_i\times\F,
\qquad K_i=D_i\sqcup\{\infty_i\},
\qquad K=\prod_{i=0}^{r}K_i.
\]

Give \(K_i\) the one-point compactification topology: every point of
\(D_i\) is isolated, and the neighborhoods of \(\infty_i\) contain
all but finitely many points of \(D_i\). The spaces \(K_i\) and their
finite product \(K\) are profinite.

Let \(\tau\) flip the bit in every finite coordinate and fix each
point at infinity:

\[
\tau(a,s)=(a,s+1),\qquad \tau(\infty_i)=\infty_i.
\]

This defines a continuous involution of \(K\). Its only fixed point
is \(p=(\infty_0,\ldots,\infty_r)\), because any finite coordinate
changes its bit. Write

\[
q:K\longrightarrow\overline K=K/\langle\tau\rangle,
\qquad U_r=\overline K\setminus\{q(p)\}.
\]

Taking the compact quotient first makes the topology of the example
particularly transparent. The quotient is compact Hausdorff: its orbit
relation is the union of the diagonal and the graph of a homeomorphism,
hence is closed. It is also zero-dimensional. Indeed, in the inverse
image of any open neighborhood of an orbit, choose a clopen neighborhood
of one representative; its union with its translate is invariant and
clopen, and descends to a clopen neighborhood in the quotient. The
orbit map is open, since the saturation of an open set \(O\) is
\(O\cup\tau(O)\). Thus \(\overline K\) is profinite and \(U_r\)
is locally profinite Hausdorff.\leanmark{construction}

\subsection{Choosing a representative on each chart}

Let \(W_i\subseteq U_r\) consist of the orbits whose \(i\)-th
coordinate is finite. These opens cover \(U_r\): the only orbit with
no finite coordinate was removed. On \(W_i\), each orbit has exactly
one representative whose \(i\)-th bit is zero. This representative
depends continuously on the orbit, because the bit-zero slice is open
and the quotient map restricts to an open continuous bijection from
that slice onto \(W_i\). We call this the \emph{gauge choice} on
\(W_i\).\leanmark{charts}

\begin{lemma}\label{lem:cells}
For every nonempty \(I\subseteq\{0,\ldots,r\}\), the intersection
\(W_I=\bigcap_{i\in I}W_i\) is a disjoint union of compact-open
subsets of \(U_r\).
\end{lemma}

\begin{proof}
Choose an anchor \(a\in I\) and use its gauge choice. Prescribe every
finite coordinate \((a_i,s_i)\in D_i\) for \(i\in I\), with
\(s_a=0\), and leave all coordinates outside \(I\) unrestricted.
This gives a clopen cylinder in \(K\). Its image in the quotient is
compact and open, and avoids \(q(p)\).

These images are disjoint for distinct prescriptions. If two meet,
the corresponding representatives either agree or are exchanged by
\(\tau\). Exchange is impossible because both anchor bits are zero;
agreement forces identical prescriptions. Every orbit in \(W_I\)
has one such prescription, so the images cover \(W_I\).
\end{proof}

These compact-open pieces explain why the cover will compute
cohomology. The intersections need not themselves be compact; it is
their decomposition into disjoint compact pieces that matters.

\subsection{Computing cardinality and weight}

\begin{proposition}\label{prop:size}
The space \(U_r\) satisfies \(|U_r|=w(U_r)=\aleph_r\).
\end{proposition}

\begin{proof}
The finite product \(K\) has cardinality \(\aleph_r\), the size of its
largest factor. A quotient and then a subspace cannot increase
cardinality, so \(|U_r|\leq\aleph_r\).

For the weight estimate, recall that an infinite compact Hausdorff
space \(T\) has \(w(T)\leq|T|\). Here is a useful direct proof.
For each ordered pair of distinct points \(x,y\), choose disjoint
neighborhoods \(O_{xy}\) of \(x\) and \(V_{xy}\) of \(y\).
Given an open neighborhood \(N\) of \(x\), finitely many of the \(V_{xy}\)
cover the compact set \(T\setminus N\). The intersection of the
corresponding \(O_{xy}\) is contained in \(N\). These finite
intersections, for all \(x\), form a basis with at most \(|T|\)
elements.

An open quotient map does not increase weight, because the images of
basis elements form a basis. Subspaces do not increase weight either.
Consequently,

\[
w(U_r)\leq w(\overline K)\leq w(K)\leq\aleph_r.
\]

For both reverse inequalities, consider the axis whose only finite
coordinate is the last one. Every orbit there has a unique
representative with last coordinate \((a,0)\), so the axis is a copy
of \(A_r\). It is a discrete subspace: the open condition that the
last base coordinate equal \(a\) isolates its corresponding axis
point. We obtain \(\aleph_r\) distinct points, and any basis for
\(U_r\) must contain distinct neighborhoods isolating these points
in the axis. Thus both cardinality and weight are at least
\(\aleph_r\).\leanmark{size}
\end{proof}
